\section{Demo 实施计划与交付物}
文档层面，本次交付首先产出一个自包含的 XeLaTeX 文档包，用于固定 proposal 结构、图示位置、里程碑和验证标准；代码层面，则为后续 notebook / script demo 预留目录与输出接口。建议的仓库交付物分为两层：
\begin{itemize}
  \item \textbf{当前已生成的文档包}
  \begin{itemize}
    \item \texttt{latex/}
    \item 根目录 \texttt{README.md}
    \item \texttt{docs/Demo\_Documentation\_Build\_Report.md}
  \end{itemize}
  \item \textbf{后续 demo 运行产物}
  \begin{itemize}
    \item \texttt{outputs/topk\_candidates.csv}
    \item \texttt{outputs/spectrum\_match.png}
    \item \texttt{outputs/molecule\_grid.png}
    \item \texttt{outputs/trace\_log.json}
  \end{itemize}
\end{itemize}

\begin{keybox}[建议目录视图]
\begin{itemize}
  \item \texttt{latex/main.tex}
  \item \texttt{latex/sections/}
  \item \texttt{latex/figures/}
  \item \texttt{latex/tables/}
  \item \texttt{latex/refs.bib}
  \item \texttt{latex/Makefile}
\end{itemize}
\end{keybox}

建议将开发过程按里程碑推进：先稳定文档与接口，再补齐最小数据流、编码器 forward、候选生成与重排，可视化与 trace log 最后闭环。对应里程碑见表~\ref{tab:milestones}。

\begin{table}[htbp]
  \centering
  \caption{Demo 里程碑规划}
  \label{tab:milestones}
  \small
  \renewcommand{\arraystretch}{1.2}
  \begin{tabular}{>{\raggedright\arraybackslash}p{0.10\linewidth} >{\raggedright\arraybackslash}p{0.22\linewidth} >{\raggedright\arraybackslash}p{0.28\linewidth} >{\raggedright\arraybackslash}p{0.28\linewidth}}
    \toprule
    ID & 里程碑 & 目标 & 预期交付物 \\
    \midrule
    M1 & Document package & 固定 proposal 结构、图注、风险边界与构建方式 & LaTeX 文档包、README、build report \\
    M2 & Minimal data pipeline & 统一输入谱图格式、样例数据与预处理接口 & 输入样例、预处理脚本或 notebook 单元 \\
    M3 & Spectrum encoder forward pass & 跑通 coordinate-aware encoder 的前向流程 & 编码张量、trace log、模块说明 \\
    M4 & Candidate retrieval / generation & 建立候选召回与基础化学过滤闭环 & 候选列表、合法性过滤结果 \\
    M5 & Forward spectrum reranking & 使用近似前向谱图或模板完成 Top-K 重排 & Top-K csv、相似度分数、对比图 \\
    M6 & Visualization and trace log & 输出结果面板与可追踪日志 & 分子网格图、谱图叠图、trace\_log.json \\
    M7 & Optional graph diffusion prototype & 预留未来 insert/delete graph diffusion 原型接口 & 原型说明、伪代码或最小实验记录 \\
    \bottomrule
  \end{tabular}
\end{table}


此外，文档中所有图均先以占位框表示，方便后续替换为正式架构图、流程图和结果图。建议在实际补图时统一采用英文图内标签、中文正文解释的形式，保证技术感与国际化可读性兼顾。

\begin{figure}[htbp]
  \centering
  \placeholderfigure{Figure placeholder: expected demo outputs}
  \caption{Expected demo outputs: input spectrum, Top-K molecular candidates, spectrum overlay, and trace log.}
\end{figure}
